\chapter{Introducción}
\label{chap:introduccion}

\section{Motivación}

Los asistentes conversacionales basados en modelos de lenguaje de gran escala
(LLM) han evolucionado desde sistemas centrados en generar texto hasta agentes
capaces de mantener estado, planificar y utilizar herramientas. Esta evolución
permite plantear interfaces en lenguaje natural para procesos administrativos
que tradicionalmente requieren navegar por formularios, conocer la organización
de un servicio o recurrir a atención telefónica. La gestión de citas sanitarias
es un caso especialmente representativo: la operación parece sencilla, pero
combina diversidad de usuarios, información sensible, restricciones de agenda,
integración con sistemas heterogéneos y consecuencias reales si una acción se
ejecuta de forma incorrecta.

Un LLM aislado no ofrece por sí mismo las garantías necesarias. Puede interpretar
mal una fecha, mantener una preferencia obsoleta, presentar como disponible un
hueco inexistente o redactar una confirmación aunque la operación haya fallado.
Además, una respuesta bien escrita puede revelar datos a alguien que no tiene
permiso para consultarlos o utilizar información para un fin no autorizado.
Por ello, el interés de este trabajo no reside únicamente en
mejorar la conversación, sino en estudiar una arquitectura que separe
interpretación, personalización, validación y ejecución.

La literatura sobre agentes basados en LLM distingue componentes de percepción,
razonamiento, memoria y acción, y muestra que la colaboración entre roles puede
incorporar capacidades complementarias \citep{xi2023agents,wu2024autogen}. Sin
embargo, añadir agentes también introduce coordinación, latencia y nuevas rutas
de fallo. En paralelo, la personalización mediante recuperación de información
del usuario mejora diversas tareas, pero su utilidad depende de seleccionar
datos relevantes y no de volcar un perfil completo en el contexto
\citep{salemi2024lamp}. Esto plantea la siguiente pregunta:
\textquestiondown{}qué ventajas y dificultades tiene usar varios agentes, en lugar
de uno solo, para gestionar citas de forma personalizada y segura?

\section{Problema y preguntas de investigación}

El trabajo estudia cómo diseñar y evaluar un asistente que gestione citas
sanitarias y adapte la conversación a las preferencias del usuario. Para
completar una gestión, el sistema debe interpretar la petición, utilizar la
información permitida y ejecutar correctamente las operaciones sobre la agenda.

Se comparan distintas arquitecturas en un entorno simulado, con perfiles y
agendas sintéticos. La evaluación busca responder las siguientes preguntas:

\begin{description}[style=nextline]
  \item[PI1.] \textquestiondown{}Qué cambia al utilizar cinco agentes especializados frente a un
  único agente, cuando ambos disponen del mismo modelo, herramientas y controles?
  \item[PI2.] \textquestiondown{}Se tienen en cuenta las preferencias del usuario cuando
  corresponde y se descartan cuando no deben utilizarse?
  \item[PI3.] \textquestiondown{}Los controles del sistema evitan accesos y operaciones no
  autorizados sin impedir gestiones legítimas?
  \item[PI4.] \textquestiondown{}Cómo evaluar conjuntamente si la gestión se completa, la
  respuesta se adapta al usuario y se respetan las restricciones del sistema?
\end{description}

\section{Objetivos}

Para responder a estas preguntas, se plantean los siguientes objetivos:

\begin{enumerate}
  \item Modelar el dominio de citas y los flujos de solicitud, consulta,
        reprogramación, cancelación y recordatorio.
  \item Convertir los requisitos relevantes del RGPD, del Reglamento de IA y de
        la Ley 41/2002 en reglas técnicas y criterios de prueba.
  \item Definir responsabilidades, contratos y rutas de escalado para los
        componentes de contexto, perfil, cumplimiento, personalización y
        orquestación.
  \item Implementar un simulador con datos sintéticos y un subconjunto de
        recursos de FHIR R4 (\textit{Fast Healthcare Interoperability
        Resources}), estándar de HL7 para intercambiar información sanitaria
        \citep{fhirR4}.
  \item Construir tres variantes comparables: una arquitectura con cinco
        agentes LLM especializados (M5), una con un único agente LLM (S1) y otra
        híbrida, con un flujo programado y tareas LLM acotadas (H).
  \item Evaluar éxito extremo a extremo, seguridad funcional, cumplimiento,
        calidad de personalización, tiempo de inferencia y consumo.
  \item Facilitar la reproducción de las pruebas mediante el código, los datos
        sintéticos y la configuración utilizada.
\end{enumerate}

\section{Alcance}

El alcance incluye identificación y representación simuladas, consulta de
disponibilidad y citas propias, reserva, reprogramación, cancelación y
recordatorios sintéticos. Las preferencias pueden ordenar huecos compatibles
o adaptar la respuesta. FHIR se utiliza para representar las citas y la agenda
con estructuras definidas para el intercambio de información sanitaria, en
lugar de crear formatos propios. El simulador implementa solo los recursos y
campos necesarios para las gestiones estudiadas. El escalado detiene el flujo
y orienta hacia atención humana, sin una cola conectada con profesionales.

La implementación disponible ejecuta las cinco gestiones conversacionales en
las agendas sintéticas de medicina de familia, pediatría y enfermería, que solo
ofrecen citas presenciales. El sistema reconoce las peticiones de citas por
teléfono o videollamada, pero estas modalidades no tienen huecos disponibles en
esas agendas. Los recordatorios se configuran en el simulador, sin enviar
mensajes externos. La consulta autorizada es de lectura. Las operaciones con
efectos requieren confirmar una propuesta válida.

El asistente no realiza diagnósticos ni pronósticos y tampoco prescribe
tratamientos. No decide la urgencia de una situación, la prioridad de un
paciente o su acceso a una prestación sanitaria. Trabaja con datos sintéticos y
no se conecta a sistemas sanitarios reales. Su función es apoyar la gestión de
citas sin sustituir la atención de profesionales sanitarios o administrativos.
Si detecta una posible emergencia, detiene la gestión e indica los canales
oficiales de atención definidos para ese caso.

\section{Metodología de desarrollo y evaluación}

El desarrollo comenzó con el análisis de las gestiones de citas, la normativa
y los trabajos relacionados. A partir de esos requisitos se construyó una
primera versión funcional. Después se incorporaron los perfiles, los agentes,
los controles y la interfaz CitaLab. Las pruebas permitieron detectar y
corregir fallos durante la integración.

Para la evaluación se elaboró un corpus sintético propio con casos en español
e inglés. Se comparan tres arquitecturas (M5, S1 y H) con dos modelos:
GPT-OSS~20B \citep{openai2025gptoss}, utilizado a través de Groq, y
Gemini~3.5 Flash-Lite \citep{gemini35FlashLiteDocs}, a través de Google.
Las variantes resuelven los mismos casos con el mismo perfil, agenda,
permisos y herramientas. Dentro de cada modelo, esas condiciones permiten
comparar las arquitecturas.

Un usuario simulado genera los mensajes a partir de un objetivo y de la
conversación que puede ver. Responde a lo que pregunta el asistente sin conocer
sus decisiones internas. Antes de la comparación, las pruebas y el piloto
comprueban el simulador y los instrumentos de evaluación. Las versiones de
código, datos y \emph{prompts} quedan registradas para cada ejecución.

Las comprobaciones automáticas verifican las operaciones realizadas y el
estado de la agenda. Un jurado de dos LLM valora la conversación, su
adaptación al usuario y si el asistente comunica lo que realmente ha ocurrido,
y una revisión humana ciega de una muestra comprueba sus notas. También
se analizan las llamadas, los tokens y el tiempo empleado. Los errores y las
interrupciones se conservan en los resultados. La
Tabla~\ref{tab:trazabilidad-preguntas} resume qué se comprueba para cada
pregunta. El capítulo~\ref{chap:metodologia-evaluacion} explica el procedimiento
completo y sus límites.

\begin{table}[H]
\centering
\caption{Trazabilidad entre preguntas y evidencia experimental.}
\label{tab:trazabilidad-preguntas}
\small
\begin{tabularx}{\textwidth}{L{0.09\textwidth}Y}
\toprule
\textbf{PI} & \textbf{Qué se comprueba} \\
\midrule
PI1 & Resultado de las gestiones, errores y recursos utilizados por cada arquitectura dentro de un mismo modelo. \\
PI2 & Aplicación de las preferencias cuando corresponden y descarte cuando no deben utilizarse. \\
PI3 & Accesos u operaciones no autorizados y bloqueos injustificados de gestiones legítimas. \\
PI4 & Correspondencia entre conversación, operaciones y agenda, junto con la revisión de la personalización y las restricciones. \\
\bottomrule
\end{tabularx}
\end{table}

\section{Contribución}

La contribución defendible no es la mera combinación de agentes, preferencias y
FHIR, ya presente de forma parcial en trabajos recientes. Se propone integrar y
evaluar, en un mismo prototipo, una combinación que no se encontró en la
revisión realizada: responsabilidades especializadas, memoria personal
selectiva y vigente, consulta y reserva de citas mediante operaciones
predefinidas y controladas por el sistema,
controles derivados de normativa europea y una traza mínima que permita atribuir
propuesta, validación y ejecución sin copiar contenido personal innecesario.

La aportación es técnica y experimental. Además del prototipo, incluye un
procedimiento de evaluación que contrasta entre sí el estado final de la
agenda, las trazas de cada decisión, un jurado de dos LLM y una revisión humana
ciega. El análisis se limita al corpus
sintético y la configuración evaluada. No demuestra preparación clínica ni
superioridad general. Permite comparar resultado, coste de coordinación y
fallos observados antes de una validación con organizaciones o usuarios reales.

\section{Resultados principales}
\label{sec:intro-resultados}

La campaña final ejecutó 276 episodios, 46 por cada combinación de arquitectura
y modelo, con un corpus distinto del que se usó para ajustar las instrucciones.
Para abreviar, se llama Groq a la configuración con GPT-OSS~20B y Gemini a la de
Gemini~3.5 Flash-Lite. Los resultados responden a las preguntas de
investigación dentro del entorno sintético estudiado, y la
Tabla~\ref{tab:respuestas-pi} los resume junto con sus limitaciones.

\paragraph{PI1. Agentes especializados.}
M5 no obtiene una ventaja significativa en el éxito de las gestiones,
comprobado sobre la agenda. Con Groq, M5 y H completan correctamente 43 de los
46 episodios y S1 completa 42. Con Gemini, M5 y H completan 43 y S1 los 46. En
cambio, M5 hace el doble de llamadas que S1 y H, usa unas 2,4 veces los tokens
de H y tiene el mayor tiempo de inferencia. Al reunir los dos modelos, H
obtiene además notas iguales o superiores a las de M5 en las seis dimensiones
de calidad, tanto en la revisión humana como en el jurado. En estas gestiones, repartir el trabajo entre cinco agentes no compensa
su coste. La explicación más probable es que los fallos más frecuentes se
originan al interpretar la petición y al actualizar la tarea, dos pasos que M5
comparte con H, y que los agentes que añade M5 actúan después de ellos.

\paragraph{PI2. Preferencias.}
Cuando la petición dejaba la hora al perfil, la cita final quedó en la franja
guardada en los 15 episodios de Gemini y en 8 de los 15 de Groq. Con Groq, M5 y
H ordenaron bien los huecos, pero la redacción quitó en ocho de diez mensajes la
frase que explicaba ese orden, y sin esa explicación el usuario simulado, que no
conocía su preferencia, tendió a elegir la opción más temprana. S1 con Groq,
además, solo consultó el perfil en tres de los siete episodios en los que hacía
falta. Cuando la petición contradecía el perfil, el asistente siguió la
petición en 11 de los 12 episodios. Para personalizar no basta con ordenar bien
las opciones, también hay que explicar por qué se ofrece primero una de ellas.

\paragraph{PI3. Controles.}
Las comprobaciones de seguridad del oráculo no detectaron ningún incumplimiento
en los 276 episodios: ninguna escritura se hizo sin confirmación, las citas de
otras personas no cambiaron y no hubo cambios en otras citas propias. Tampoco se
bloqueó ninguna gestión legítima. Esas comprobaciones no cubren todo el
contenido de la respuesta, y en un episodio el asistente añadió consejos
clínicos que el oráculo no detectó. Además, una confirmación correcta no
garantiza que se cumpla la intención de la persona. En nueve episodios, una
petición ambigua para mover una cita se convirtió en una reserva nueva,
propuesta con claridad, aceptada y autorizada. La comprobación que revisa la
redacción antes de mostrarla también protege menos de lo que parece. De los 32
borradores que sustituyó por parecer anuncios de un efecto, solo tres
anunciaban o prometían algo que no había ocurrido.

\paragraph{PI4. Evaluación conjunta.}
Ningún instrumento bastó por sí solo, porque cada uno detectó problemas que los
demás no veían. El oráculo detectó los 16 episodios fallidos, aunque en 11 de
ellos el usuario simulado había dado por conseguido su objetivo. El jurado de
dos LLM separó casi por completo los fallos de los éxitos y detectó un falso
éxito del oráculo. Su concordancia con la revisión humana es moderada o buena,
con un acuerdo exacto de entre el 72 y el 95\,\% según el juez y la dimensión,
y el jurado es algo más severo. Las trazas permitieron localizar el origen de
los errores, repartido entre la interpretación de la petición, la
actualización del estado, los contratos de las herramientas, las instrucciones
que construye el código y la redacción.

\paragraph{Alcance.}
Estos resultados se limitan al corpus sintético, a los dos modelos y a las
gestiones estudiadas. No permiten establecer una clasificación general de las
arquitecturas ni sacar conclusiones sobre pacientes reales, la seguridad
clínica o el cumplimiento jurídico de un despliegue. El
capítulo~\ref{chap:conclusiones} desarrolla estas conclusiones, junto con las
limitaciones y el trabajo futuro.

\section{Estructura de la memoria}

El capítulo~\ref{chap:fundamentos} define IA agéntica y personalización.
El capítulo~\ref{chap:dominio-regulado} delimita dominio, datos y requisitos.
El capítulo~\ref{chap:estado-arte} compara antecedentes y deriva criterios de
diseño. Los capítulos~\ref{chap:diseno} y~\ref{chap:implementacion} justifican
la arquitectura y explican su realización, respectivamente. El
capítulo~\ref{chap:metodologia-evaluacion} describe los instrumentos y cálculos.
El capítulo~\ref{chap:resultados} presenta la cobertura de la evaluación, los
resultados por configuración, el análisis de errores y un estudio de casos. El
capítulo~\ref{chap:discusion} interpreta esos resultados, propone y comprueba
correcciones, responde a las preguntas de investigación y examina las amenazas
a la validez. El
capítulo~\ref{chap:conclusiones} recoge las conclusiones, el grado de
cumplimiento de los objetivos, las limitaciones y el trabajo futuro. Los
apéndices recogen el inventario normativo con la trazabilidad de requisitos y
los criterios de revisión de las conversaciones.
